\documentclass[11pt,aspectratio=169]{beamer}
\usetheme{SimplePlus}

\usepackage[T1]{fontenc}
\usepackage[utf8]{inputenc}
\usepackage{amsmath,amssymb,mathtools}
\usepackage{booktabs}
\usepackage{array}
\usepackage{appendixnumberbeamer}
\usepackage{hyperref}

\newcommand{\E}{\mathbb{E}}
\newcommand{\Var}{\mathrm{Var}}
\newcommand{\Cov}{\mathrm{Cov}}
\newcommand{\Prob}{\mathbb{P}}
\newcommand{\Corr}{\mathrm{Corr}}
\newcommand{\AR}{\mathrm{AR}}
\newcommand{\VAR}{\mathrm{VAR}}
\newcommand{\MA}{\mathrm{MA}}
\newcommand{\IRF}{\mathrm{IRF}}
\newcommand{\AIC}{\mathrm{AIC}}
\newcommand{\BIC}{\mathrm{BIC}}
\newcommand{\OLS}{\mathrm{OLS}}
\newcommand{\plim}{\operatorname*{plim}}
\newcommand{\derivbutton}[1]{\vspace{0.4em}\hfill\hyperlink{app:#1}{\beamerbutton{Appendix details}}}
\newcommand{\backbutton}[1]{\vspace{0.4em}\hfill\hyperlink{main:#1}{\beamerbutton{Back to main slide}}}

\title{Applied Econometrics: Session 10}
\subtitle{A Basic Introduction to Macroeconometrics}
\author{PhD Intensive Course}
\institute{Lecture 10}
\date{\today}

\begin{document}

\begin{frame}
  \titlepage
\end{frame}

\begin{frame}{How This Lecture Fits the Course}
  The first part of the course focused on micro-style causal designs: experiments, controls, IV, matching, DiD, synthetic control, and RD.
  \vspace{0.5em}

  Macroeconometrics asks related questions with different data:
  \begin{itemize}
    \item observations are usually indexed by time: \(t=1,\ldots,T\);
    \item outcomes are persistent and move together;
    \item one aggregate shock can affect many variables at once;
    \item sample sizes are often small: 60 quarterly observations is only 15 years.
  \end{itemize}

  \begin{block}{Core message}
    In macroeconometrics, time dependence is not a nuisance. It is usually the object of interest.
  \end{block}
\end{frame}

\begin{frame}{Roadmap for a 2.5 Hour Introduction}
  \begin{tabular}{p{0.23\linewidth}p{0.68\linewidth}}
    \toprule
    Time & Topic \\
    \midrule
    0--20 min & What macroeconometrics studies; why time-series data are special \\
    20--55 min & Trends, growth rates, stationarity, autocorrelation, and unit roots \\
    55--90 min & Univariate dynamics: AR models, persistence, and forecasts \\
    90--125 min & Multivariate dynamics: VARs, Granger causality, and impulse responses \\
    125--145 min & Identification: why a VAR residual is not automatically a structural shock \\
    145--150 min & Checklist and takeaways \\
    \bottomrule
  \end{tabular}
\end{frame}

\begin{frame}{Learning Goals}
  By the end, students should be able to:
  \begin{itemize}
    \item distinguish forecasting, description, and causal/structural analysis;
    \item explain why serial correlation changes standard econometric reasoning;
    \item recognize trends, growth rates, stationarity, and unit-root behavior;
    \item estimate and interpret a simple autoregression;
    \item explain what a VAR is and what an impulse response measures;
    \item state why macro shocks require identifying assumptions.
  \end{itemize}
\end{frame}

\section{What Is Macroeconometrics?}

\begin{frame}{What Is Macroeconometrics?}
  Macroeconometrics is econometrics for aggregate time-series and panel-time-series data.

  \begin{itemize}
    \item GDP, consumption, investment, inflation, unemployment, wages.
    \item Interest rates, credit, exchange rates, house prices, asset prices.
    \item Fiscal variables: spending, taxes, debt, deficits.
    \item Cross-country or regional macro panels.
  \end{itemize}

  \begin{block}{Simple definition}
    Macroeconometrics studies how aggregate economic variables move over time, forecast each other, and respond to shocks and policy changes.
  \end{block}
\end{frame}

\begin{frame}{Three Different Tasks}
  The same dataset can be used for three different purposes.
  \vspace{0.4em}

  \begin{tabular}{p{0.24\linewidth}p{0.34\linewidth}p{0.31\linewidth}}
    \toprule
    Task & Question & Example \\
    \midrule
    Description & What happened? & Inflation rose before unemployment fell. \\
    Forecasting & What will happen next? & Predict next quarter's GDP growth. \\
    Structural analysis & What caused what? & Estimate the effect of a monetary-policy shock. \\
    \bottomrule
  \end{tabular}

  \begin{block}{Warning}
    A good forecasting model is not automatically a good causal model.
  \end{block}
\end{frame}

\begin{frame}{Why Macro Data Are Hard}
  Macro data are not just large versions of cross-sectional data.

  \begin{itemize}
    \item \textbf{Persistence:} high GDP today predicts high GDP tomorrow.
    \item \textbf{Common shocks:} recessions move many variables together.
    \item \textbf{Policy feedback:} central banks react to inflation and output.
    \item \textbf{Small samples:} annual and quarterly data give limited \(T\).
    \item \textbf{Nonstationarity:} many levels trend over time.
    \item \textbf{Expectations:} decisions depend on anticipated future policy.
  \end{itemize}
\end{frame}

\begin{frame}{Running Example}
  We will use the same language for several macro variables:
  \[
    y_t = \text{output growth},\qquad
    \pi_t = \text{inflation},\qquad
    i_t = \text{policy interest rate}.
  \]

  Typical questions:
  \begin{itemize}
    \item Does inflation help forecast future interest rates?
    \item How persistent is output growth?
    \item What happens to output and inflation after a monetary tightening?
    \item Did two trending variables move together for economic reasons or just because both trended?
  \end{itemize}
\end{frame}

\section{Time-Series Basics}

\begin{frame}{From Cross Sections to Time Series}
  In a cross section, observations are indexed by units:
  \[
    i=1,\ldots,n.
  \]
  In a time series, observations are indexed by time:
  \[
    t=1,\ldots,T.
  \]

  \begin{block}{Key difference}
    \(y_t\) and \(y_{t-1}\) are usually related. Treating them as independent observations is often wrong.
  \end{block}
\end{frame}

\begin{frame}{Lag Notation}
  A lag is a past value of the same variable:
  \[
    y_{t-1},\quad y_{t-2},\quad y_{t-3}.
  \]

  The lag operator \(L\) is shorthand:
  \[
    Ly_t = y_{t-1},\qquad L^2y_t=y_{t-2}.
  \]

  Differences are written as
  \[
    \Delta y_t = y_t-y_{t-1}=(1-L)y_t.
  \]
\end{frame}

\begin{frame}{Always Plot the Series First}
  Before estimating a model, plot the data.

  \begin{itemize}
    \item Is the variable a level, log level, growth rate, or change?
    \item Is there a trend?
    \item Are there breaks: crises, pandemics, wars, policy-regime changes?
    \item Are there outliers or measurement changes?
    \item Is the frequency annual, quarterly, monthly, or daily?
  \end{itemize}

  \begin{block}{Practical rule}
    A macro regression that has not been preceded by a time-series plot is easy to misunderstand.
  \end{block}
\end{frame}

\begin{frame}{Levels, Logs, and Growth Rates}
  For a positive variable \(Y_t\), define the log level:
  \[
    y_t=\log(Y_t).
  \]
  The log growth rate is
  \[
    \Delta y_t
    =
    \log(Y_t)-\log(Y_{t-1})
    \approx
    \frac{Y_t-Y_{t-1}}{Y_{t-1}}.
  \]

  \begin{itemize}
    \item GDP level is usually trending.
    \item GDP growth is often closer to stationary.
    \item Inflation is already a rate of price growth.
  \end{itemize}
\end{frame}

\begin{frame}{Stationarity}
  \hypertarget{main:stationarity}{}
  A time series is weakly stationary if:
  \[
    \E[y_t]=\mu \quad \text{for all }t,
  \]
  \[
    \Var(y_t)=\sigma_y^2 \quad \text{for all }t,
  \]
  and
  \[
    \Cov(y_t,y_{t-k})
  \]
  depends on the lag \(k\), not on the calendar date \(t\).

  \begin{block}{Intuition}
    The process can be persistent, but its basic probabilistic behavior is stable over time.
  \end{block}
  \derivbutton{ar-stationarity}
\end{frame}

\begin{frame}{Autocorrelation}
  The lag-\(k\) autocorrelation is
  \[
    \rho_k
    =
    \Corr(y_t,y_{t-k})
    =
    \frac{\Cov(y_t,y_{t-k})}{\Var(y_t)}.
  \]

  \begin{itemize}
    \item If \(\rho_1\) is large, the series is persistent.
    \item If \(\rho_k\) decays slowly, shocks fade slowly.
    \item The autocorrelation function is a basic descriptive tool.
  \end{itemize}
\end{frame}

\begin{frame}{White Noise}
  A useful benchmark is white noise:
  \[
    \varepsilon_t \sim (0,\sigma^2),\qquad
    \Cov(\varepsilon_t,\varepsilon_s)=0 \quad \text{for }t\neq s.
  \]

  \begin{itemize}
    \item White noise is unpredictable from its own past.
    \item It can still have large variance.
    \item Many macro models describe observed variables as responses to shocks.
  \end{itemize}

  \begin{block}{Do not confuse}
    White noise means no serial correlation. It does not mean economically unimportant.
  \end{block}
\end{frame}

\section{Univariate Dynamics}

\begin{frame}{The AR(1) Model}
  \hypertarget{main:ar1}{}
  The simplest persistent time-series model is
  \[
    y_t = c+\phi y_{t-1}+\varepsilon_t,
    \qquad \varepsilon_t \sim (0,\sigma^2).
  \]

  \begin{itemize}
    \item \(c\) controls the unconditional mean.
    \item \(\phi\) measures persistence.
    \item \(\varepsilon_t\) is the one-period-ahead forecast error.
  \end{itemize}

  \begin{block}{Stationarity condition}
    The AR(1) is stationary if \(|\phi|<1\).
  \end{block}
  \derivbutton{ar-stationarity}
\end{frame}

\begin{frame}{Interpreting Persistence}
  Suppose
  \[
    y_t=c+\phi y_{t-1}+\varepsilon_t.
  \]

  If a shock raises \(y_t\) by 1 today, the expected effect is:
  \[
    1,\quad \phi,\quad \phi^2,\quad \phi^3,\ldots
  \]
  over future periods.

  \begin{itemize}
    \item \(\phi=0.2\): shock disappears quickly.
    \item \(\phi=0.8\): shock is persistent.
    \item \(\phi=1\): shock never dies out; this is a unit root.
  \end{itemize}
  \derivbutton{ar-forecast}
\end{frame}

\begin{frame}{Forecasting with an AR(1)}
  The one-step-ahead forecast is
  \[
    \widehat{y}_{t+1|t}=\widehat{c}+\widehat{\phi}y_t.
  \]

  The \(h\)-step-ahead forecast is
  \[
    \widehat{y}_{t+h|t}
    =
    \widehat{\mu}
    +
    \widehat{\phi}^{\,h}(y_t-\widehat{\mu}),
    \qquad
    \widehat{\mu}=\frac{\widehat{c}}{1-\widehat{\phi}}.
  \]

  \begin{block}{Interpretation}
    Forecasts revert to the mean if \(|\widehat{\phi}|<1\). They do not revert if \(\widehat{\phi}\) is close to one.
  \end{block}
\end{frame}

\begin{frame}{Estimating the AR(1)}
  Estimate
  \[
    y_t=c+\phi y_{t-1}+\varepsilon_t
  \]
  by OLS on observations \(t=2,\ldots,T\).

  \begin{itemize}
    \item Dependent variable: \(y_t\).
    \item Regressor: \(y_{t-1}\).
    \item Coefficient: persistence.
    \item Residual: forecast error, not necessarily a structural shock.
  \end{itemize}

  \begin{block}{Small-sample warning}
    With highly persistent data and small \(T\), estimates and standard errors can be fragile.
  \end{block}
\end{frame}

\begin{frame}{Unit Roots}
  A unit-root process has
  \[
    y_t=y_{t-1}+\varepsilon_t.
  \]

  Then
  \[
    y_t=y_0+\sum_{s=1}^{t}\varepsilon_s.
  \]

  \begin{itemize}
    \item Shocks have permanent effects.
    \item The variance grows with time.
    \item Usual stationary regression intuition breaks down.
  \end{itemize}
\end{frame}

\begin{frame}{Spurious Regression}
  \hypertarget{main:spurious}{}
  If two unrelated variables both trend or both have unit roots, the regression
  \[
    y_t=\alpha+\beta x_t+u_t
  \]
  can look impressive even when \(x_t\) has no economic effect on \(y_t\).

  \begin{block}{Symptom}
    High \(R^2\), significant \(t\)-statistics, and strongly autocorrelated residuals in a regression of unrelated trending series.
  \end{block}

  \begin{itemize}
    \item Plot the data.
    \item Consider growth rates or first differences.
    \item Check whether there is a cointegrating relationship.
  \end{itemize}
  \derivbutton{spurious}
\end{frame}

\begin{frame}{Cointegration}
  Two series can each be nonstationary, but a linear combination can be stationary:
  \[
    y_t-\beta x_t \quad \text{is stationary}.
  \]

  \begin{block}{Intuition}
    The two variables may drift over time, but they do not drift arbitrarily far apart.
  \end{block}

  Examples where cointegration may be plausible:
  \begin{itemize}
    \item consumption and income;
    \item short and long interest rates;
    \item prices in closely connected markets;
    \item money, prices, and output under some monetary theories.
  \end{itemize}
\end{frame}

\section{Multivariate Dynamics}

\begin{frame}{Why Multivariate Models?}
  Macro variables are jointly determined.

  \begin{itemize}
    \item Inflation affects monetary policy.
    \item Monetary policy affects output and inflation.
    \item Output affects unemployment, tax revenue, and credit.
    \item Financial variables react quickly to expectations.
  \end{itemize}

  \begin{block}{Problem}
    A single-equation model can hide feedback among variables.
  \end{block}
\end{frame}

\begin{frame}{The VAR(1) Model}
  \hypertarget{main:var1}{}
  Let
  \[
    z_t=
    \begin{bmatrix}
      y_t\\
      \pi_t\\
      i_t
    \end{bmatrix}.
  \]

  A VAR(1) is
  \[
    z_t=a+A z_{t-1}+u_t,
    \qquad
    \E[u_t]=0,\quad
    \E[u_tu_t']=\Sigma_u.
  \]

  \begin{itemize}
    \item Each variable is regressed on lagged values of all variables.
    \item The matrix \(A\) summarizes dynamic interactions.
    \item \(u_t\) is a vector of reduced-form forecast errors.
  \end{itemize}
  \derivbutton{var-companion}
\end{frame}

\begin{frame}{VAR Equations in Words}
  A three-variable VAR says:
  \[
    \text{output today}
    =
    \text{lags of output, inflation, interest rates}
    +
    \text{forecast error}.
  \]
  \[
    \text{inflation today}
    =
    \text{lags of output, inflation, interest rates}
    +
    \text{forecast error}.
  \]
  \[
    \text{interest rate today}
    =
    \text{lags of output, inflation, interest rates}
    +
    \text{forecast error}.
  \]

  \begin{block}{Beginner intuition}
    A VAR is a disciplined way to let macro variables predict each other dynamically.
  \end{block}
\end{frame}

\begin{frame}{Choosing the Lag Length}
  A VAR(\(p\)) uses \(p\) lags:
  \[
    z_t=a+A_1z_{t-1}+\cdots+A_pz_{t-p}+u_t.
  \]

  More lags can capture richer dynamics but cost degrees of freedom.

  \begin{itemize}
    \item Use economic judgment about timing.
    \item Check residual autocorrelation.
    \item Compare information criteria such as \(\AIC\) and \(\BIC\).
    \item Remember that macro samples are often small.
  \end{itemize}
\end{frame}

\begin{frame}{Granger Causality}
  Variable \(x_t\) Granger-causes \(y_t\) if lags of \(x_t\) help forecast \(y_t\) after controlling for lags of \(y_t\).

  Example:
  \[
    y_t=\alpha+\rho y_{t-1}+\gamma x_{t-1}+u_t.
  \]
  Test
  \[
    H_0:\gamma=0.
  \]

  \begin{block}{Important}
    Granger causality is predictive content. It is not automatically causal in the treatment-effect sense.
  \end{block}
\end{frame}

\begin{frame}{Impulse Responses}
  \hypertarget{main:irf}{}
  An impulse response traces what happens after a shock.

  \begin{itemize}
    \item Shock: an unexpected movement in one variable or structural disturbance.
    \item Response: predicted path of variables after the shock.
    \item Horizon: period \(0,1,2,\ldots,H\).
  \end{itemize}

  For a VAR(1), if the initial shock is \(s\), then
  \[
    \IRF(h)=A^h s.
  \]

  \derivbutton{irf-recursion}
\end{frame}

\begin{frame}{Reduced-Form Errors Are Not Structural Shocks}
  A VAR residual \(u_t\) is a forecast error:
  \[
    z_t-\E[z_t\mid z_{t-1},z_{t-2},\ldots].
  \]

  But a structural shock is an economically meaningful disturbance:
  \begin{itemize}
    \item monetary-policy shock;
    \item technology shock;
    \item demand shock;
    \item fiscal spending shock.
  \end{itemize}

  \begin{block}{Identification problem}
    The data give reduced-form forecast errors. Turning them into structural shocks requires assumptions.
  \end{block}
\end{frame}

\begin{frame}{Common Identification Strategies}
  Macroeconomists identify shocks using restrictions or external information.

  \begin{itemize}
    \item \textbf{Recursive ordering:} some variables react within the period, others react with delay.
    \item \textbf{Sign restrictions:} a monetary tightening should raise rates and lower prices/output over some horizon.
    \item \textbf{External instruments:} use high-frequency surprises or narrative measures as shock proxies.
    \item \textbf{Natural experiments:} use institutional rules or policy discontinuities when available.
  \end{itemize}

  \begin{block}{Connection to earlier lectures}
    This is the macro version of the same causal problem: what variation identifies the effect?
  \end{block}
\end{frame}

\begin{frame}{Local Projections}
  A popular alternative to VAR impulse responses is the local projection:
  \[
    y_{t+h}=\alpha_h+\beta_h shock_t+\Gamma_h'controls_t+e_{t+h}.
  \]

  Estimate a separate regression for each horizon \(h\).

  \begin{itemize}
    \item Flexible and easy to explain.
    \item Works naturally with nonlinearities and interactions.
    \item Still requires an identified shock.
    \item Standard errors must handle serial correlation across horizons.
  \end{itemize}
\end{frame}

\section{Practical Workflow}

\begin{frame}{A Basic Macroeconometric Workflow}
  \begin{enumerate}
    \item State the question: description, forecast, or structural effect.
    \item Define variables, frequency, sample, transformations, and timing.
    \item Plot all variables.
    \item Decide levels, logs, growth rates, or differences.
    \item Estimate a simple benchmark model.
    \item Check residuals and robustness.
    \item For causal claims, state the identifying assumptions.
    \item Report magnitudes in economic units.
  \end{enumerate}
\end{frame}

\begin{frame}{Python Example for This Lecture}
  The script
  \[
    \texttt{lectures/code/lecture-10-macroeconometrics.py}
  \]
  demonstrates three ideas:
  \begin{itemize}
    \item estimating an AR(1) and forecasting;
    \item seeing a spurious regression from unrelated random walks;
    \item estimating a two-variable VAR and computing impulse responses.
  \end{itemize}

  \begin{block}{Student task}
    Change the persistence parameter and watch how forecasts and impulse responses change.
  \end{block}
\end{frame}

\begin{frame}{What to Report in a Basic Macro Exercise}
  Report enough for the reader to reconstruct the analysis:
  \begin{itemize}
    \item variables, definitions, units, and data frequency;
    \item sample start/end dates and missing-data treatment;
    \item transformations: log, difference, annualization, detrending;
    \item model specification and lag length;
    \item standard-error choice;
    \item diagnostics: residual autocorrelation, stability, outliers;
    \item identification assumptions for shocks;
    \item robustness to sample period and lag length.
  \end{itemize}
\end{frame}

\begin{frame}{Common Beginner Mistakes}
  \begin{itemize}
    \item Regressing two trending variables in levels and treating the result as causal.
    \item Calling Granger causality ``true causality.''
    \item Forgetting that policy responds to the economy.
    \item Reading a VAR residual as a structural shock without identification.
    \item Using too many lags in a small quarterly sample.
    \item Ignoring breaks such as crises or major policy-regime changes.
    \item Reporting coefficients without translating them into economic units.
  \end{itemize}
\end{frame}

\begin{frame}{Takeaways}
  \begin{enumerate}
    \item Macro data are dynamic: the past predicts the present.
    \item Stationarity matters because many standard tools rely on stable distributions.
    \item AR models teach persistence and forecasting.
    \item VARs describe joint dynamics among macro variables.
    \item Impulse responses are useful, but structural interpretation needs identification.
    \item Always separate forecasting success from causal credibility.
  \end{enumerate}
\end{frame}

\appendix

\begin{frame}{Appendix A1: AR(1) Mean and Stationarity}
  \hypertarget{app:ar-stationarity}{}
  Start with
  \[
    y_t=c+\phi y_{t-1}+\varepsilon_t,\qquad \E[\varepsilon_t]=0.
  \]
  If the mean is constant, \(\E[y_t]=\E[y_{t-1}]=\mu\). Then
  \[
    \mu=c+\phi\mu
    \quad\Rightarrow\quad
    \mu=\frac{c}{1-\phi}.
  \]
  Recursive substitution gives
  \[
    y_t=\mu+\phi^j(y_{t-j}-\mu)+\sum_{\ell=0}^{j-1}\phi^\ell\varepsilon_{t-\ell}.
  \]
  For the influence of the distant past to vanish, we need \(|\phi|<1\).
  \backbutton{ar1}
\end{frame}

\begin{frame}{Appendix A2: AR(1) Variance}
  If \(|\phi|<1\), write
  \[
    y_t-\mu=\phi(y_{t-1}-\mu)+\varepsilon_t.
  \]
  Taking variances and using \(\Cov(y_{t-1},\varepsilon_t)=0\):
  \[
    \Var(y_t)=\phi^2\Var(y_{t-1})+\sigma^2.
  \]
  Stationarity means \(\Var(y_t)=\Var(y_{t-1})=\sigma_y^2\). Hence
  \[
    \sigma_y^2=\phi^2\sigma_y^2+\sigma^2
    \quad\Rightarrow\quad
    \sigma_y^2=\frac{\sigma^2}{1-\phi^2}.
  \]
  This is finite only if \(|\phi|<1\).
  \backbutton{stationarity}
\end{frame}

\begin{frame}{Appendix A3: AR(1) Forecast Formula}
  \hypertarget{app:ar-forecast}{}
  With mean \(\mu=c/(1-\phi)\),
  \[
    y_t-\mu=\phi(y_{t-1}-\mu)+\varepsilon_t.
  \]
  The one-step conditional expectation is
  \[
    \E[y_{t+1}-\mu\mid y_t]=\phi(y_t-\mu).
  \]
  Iterating forward:
  \[
    \E[y_{t+h}\mid y_t]
    =
    \mu+\phi^h(y_t-\mu).
  \]
  The same logic gives shock responses \(1,\phi,\phi^2,\ldots\).
  \backbutton{ar1}
\end{frame}

\begin{frame}{Appendix A4: Why Unit Roots Create Trouble}
  \hypertarget{app:spurious}{}
  If
  \[
    y_t=y_{t-1}+\varepsilon_t,
  \]
  then
  \[
    y_t=y_0+\sum_{s=1}^t\varepsilon_s.
  \]
  The variance is
  \[
    \Var(y_t)=t\sigma^2
  \]
  when shocks are independent with variance \(\sigma^2\).

  Since the variance changes with \(t\), the process is not stationary. Two independent random walks can appear related because both wander persistently.
  \backbutton{spurious}
\end{frame}

\begin{frame}{Appendix A5: VAR Companion Form}
  \hypertarget{app:var-companion}{}
  A VAR(\(p\)) is
  \[
    z_t=a+A_1z_{t-1}+\cdots+A_pz_{t-p}+u_t.
  \]
  Stack current and lagged values:
  \[
    Z_t=
    \begin{bmatrix}
      z_t\\z_{t-1}\\ \vdots \\z_{t-p+1}
    \end{bmatrix}.
  \]
  Then
  \[
    Z_t=\alpha+FZ_{t-1}+U_t.
  \]
  This turns the VAR(\(p\)) into a larger VAR(1), which makes forecasting and impulse-response calculations easier.
  \backbutton{var1}
\end{frame}

\begin{frame}{Appendix A6: VAR Impulse Response Recursion}
  \hypertarget{app:irf-recursion}{}
  For a VAR(1),
  \[
    z_t=a+Az_{t-1}+u_t.
  \]
  Suppose a shock changes \(z_t\) by vector \(s\) at horizon 0.
  \[
    \IRF(0)=s.
  \]
  One period later:
  \[
    \IRF(1)=A s.
  \]
  Two periods later:
  \[
    \IRF(2)=A^2 s.
  \]
  Therefore
  \[
    \IRF(h)=A^h s.
  \]
  \backbutton{irf}
\end{frame}

\begin{frame}{Appendix A7: Cholesky Identification}
  Suppose reduced-form errors satisfy
  \[
    \E[u_tu_t']=\Sigma_u.
  \]
  Cholesky identification chooses a lower-triangular matrix \(B\) such that
  \[
    \Sigma_u=BB'.
  \]
  Then
  \[
    u_t=B\eta_t,\qquad \E[\eta_t\eta_t']=I.
  \]
  The first structural shock can affect all variables contemporaneously, the second cannot affect the first contemporaneously, and so on. This is an identifying assumption, not a fact delivered by the data.
\end{frame}

\begin{frame}{Appendix A8: Serial Correlation and Standard Errors}
  In a simple regression
  \[
    y_t=\alpha+\beta x_t+u_t,
  \]
  usual OLS standard errors assume errors are not serially correlated.

  If \(u_t\) is correlated with \(u_{t-1}\), then the usual standard error can be too small.

  Common responses:
  \begin{itemize}
    \item model the dynamics directly;
    \item use heteroskedasticity-and-autocorrelation consistent standard errors;
    \item aggregate or transform the data only when economically justified;
    \item be transparent about small-sample uncertainty.
  \end{itemize}
\end{frame}

\end{document}
